\section{Grobe Flow-Analyse}
\label{sec:flow-azn}

\subsection{Annahmen}
\label{sec:annahmen-azn}

Betr\"age sind gegeben; ihre Mechanik nie. Jede Annahme, die nicht aus den
Daten folgt, steht deshalb hier, vor der Rechnung, die sie verbraucht.

\annahme{Zwei Rechnungen, keine Summe.} Abschnitt~\ref{sec:kapazitaet-azn}
fragt, wie schnell das Unternehmen seine Schulden tilgen k\"onnte;
Abschnitt~\ref{sec:anleger-azn} fragt, was ein Anleger verdient, der heute
kauft. Beide Rechnungen verwenden \emph{denselben} freien Mittelzufluss und
sind deshalb Alternativen, keine Summanden. Wer sie addiert, gibt dasselbe
Geld zweimal aus.

\annahme{Freier Mittelzufluss.} Operativer Mittelzufluss abz\"uglich der
Investitionen in Sachanlagen und immaterielle Verm\"ogenswerte, zuz\"uglich der
Erl\"ose aus deren Abgang, abz\"uglich der Leasingtilgung. Nicht abgezogen
sind die Zahlungen auf bedingte Kaufpreise aus fr\"uheren Zuk\"aufen, die 2025
\MioUSD{\AznBedingteKaufpreise} betrugen.\quelleGB{2025}{128}\quelleMerken{s16gbxcacfxbci}
Sie sind wiederkehrend, aber sie sind Kaufpreis und nicht laufender Aufwand;
wer sie abzieht, bekommt eine um rund ein Siebtel niedrigere Reihe und
dieselbe Rangfolge.

\annahme{Konstantes Niveau, kein Wachstum.} In der Anlegerrechnung wird der
freie Mittelzufluss je Aktie \"uber den ganzen Horizont konstant gehalten. Das
ist keine Prognose, sondern der Verzicht darauf. Was der heutige Kurs an
Wachstum bereits voraussetzt, wird in Abschnitt~\ref{sec:anleger-azn}
umgekehrt ausgerechnet und gegen die tats\"achlich erzielten Raten der
Mehrjahresreihe gestellt.

\annahme{Drei Niveaus statt drei Prognosen.} Gerechnet wird mit drei
Werten, die alle aus ver\"offentlichten Zahlen stammen: den letzten zw\"olf
Monaten bis zum 30.~Juni 2026 (\USDje{\AznFcfLtm}), dem Median der sieben
Jahre (\USDje{\AznFcfMedian}) und dem letzten vollen Gesch\"aftsjahr
(\USDje{\AznFcfVoll}). Drei Zahlen, die das Unternehmen selbst berichtet hat,
muss niemand verteidigen; drei selbst gew\"ahlte Wachstumsraten schon.

\annahme{Endwert gleich Buchwert je Aktie.} Am Ende des Horizonts ist der
Anteil noch etwas wert. Angesetzt wird das bilanzielle Eigenkapital je Aktie,
\USDje{\AznBuchwertJeAktieX}. Es steht in der Bilanz, es h\"angt nicht vom
Einstiegspreis ab, und es ist konservativ; f\"ur ein Pharmaunternehmen, dessen
Verm\"ogen zu gro{\ss}en Teilen nicht bilanziert ist, sehr konservativ. Weil der
Endwert das Ergebnis beherrscht, wird die Rechnung mit einer zweiten
Endwertannahme wiederholt und beide werden ausgewiesen.

\annahme{Steuern.} Die Rechnung ist eine Betrachtung auf Unternehmensebene.
Der freie Mittelzufluss ist nach den gezahlten Ertragsteuern des Konzerns
(\MioUSD{\AznSteuerZahlung} im Jahr 2025), aber vor der Steuer des Anlegers
auf Dividende und Kursgewinn. Quellensteuern bleiben unber\"ucksichtigt, weil
der Bericht die Steuerlage des Lesers nicht kennt. Die ausgewiesenen
Renditen sind deshalb Vorsteuerrenditen \emph{des Anlegers}, und da beide
Titel gleich behandelt werden, bleibt die Rangfolge unber\"uhrt.

\annahme{Horizont.} Gerechnet wird \"uber f\"unf, zehn und f\"unfzehn Jahre ab
2026. Ein interner Zinsfu{\ss} ohne seinen Horizont ist eine Zahl, die eine
unausgesprochene Annahme trifft.

\annahme{H\"urde.} \Pct{\Huerde} pro Jahr, die Zehnjahresrendite des
Branchenindex.\quelleWeb{BlackRock, iShares U.S. Pharmaceuticals ETF, Zehnjahresrendite des Benchmark-Index}{quelle:huerde}{17.09.2026}\quelleMerken{s16webxblackrockxxisharesxuxsxxpharmaceuticalsxetfxxzehnjahresrenditexdesxbenchmarkxindexxquellexhuerde}
Sie ist eine Gesamtrendite einschlie{\ss}lich wiederangelegter
Aussch\"uttungen und damit strenger als eine reine Kursrendite.

\subsection{Entschuldungskapazit\"at}
\label{sec:kapazitaet-azn}

Die erste Frage ist eine Kapazit\"atsrechnung und keine Prognose: Wie lange
br\"auchte das Unternehmen, um seine Nettoverschuldung aus eigener Kraft zu
tilgen, wenn es die Dividende weiterzahlt und sonst nichts tut?

\begin{table}[htbp]
\centering
\caption{Entschuldungskapazit\"at, Mio.\,USD. Eine Kapazit\"atsrechnung, keine
Prognose: Sie unterstellt, dass jeder freie Dollar nach Dividende in die
Tilgung geht und keine Zuk\"aufe stattfinden.}
\label{tab:azn-kapazitaet}
\begin{tabular}{lr}
\toprule
Position & Betrag\\
\midrule
\tabellenkoerper{data/flow_azn_kapazitaet}
\bottomrule
\end{tabular}
\end{table}

\bk{Die Antwort lautet: rechnerisch unter sieben Jahre, und das ist die
uninteressantere H\"alfte des Befunds. Die interessantere ist, dass diese
Rechnung f\"ur dieses Unternehmen kaum bindet. Bei einer Verschuldung vom
\Faktor{\AznVerschuldungsgrad}-fachen des EBITDA ist die Bilanz keine
Beschr\"ankung; die Annahme \enquote{keine Zuk\"aufe} ist zugleich die
unrealistischste des Berichts, denn das Unternehmen hat 2025 mehrere
Transaktionen abgeschlossen und im Juli 2026 zwei weitere Lizenzen
gezeichnet.}

\FloatBarrier

\subsection{Die Rechnung des Anlegers}
\label{sec:anleger-azn}

Die zweite Frage ist die, die \"uber den Kauf entscheidet. Die Zahlungsreihe
lautet: Einstiegspreis heute, danach in jedem Jahr der ausschuettungsf\"ahige
Mittelzufluss je Aktie, und am Ende der Buchwert je Aktie.

\begin{table}[htbp]
\centering
\caption{Interner Zinsfu{\ss} in Prozent pro Jahr beim heutigen Kurs von
\USDje{\AznKurs}, Endwert gleich Buchwert je Aktie}
\label{tab:azn-irr}
\begin{tabular}{lrrr}
\toprule
Niveau des Mittelzuflusses & 5 Jahre & 10 Jahre & 15 Jahre\\
\midrule
\tabellenkoerper{data/flow_azn_irr}
\bottomrule
\end{tabular}
\end{table}

\begin{table}[htbp]
\centering
\caption{Einstiegsschwelle in USD: der Preis, bei dem die Rendite die H\"urde
von \Pct{\Huerde} genau erreicht. Endwert gleich Buchwert je Aktie.}
\label{tab:azn-schwelle}
\begin{tabular}{lrrrr}
\toprule
Niveau & je Aktie & 5 Jahre & 10 Jahre & 15 Jahre\\
\midrule
\tabellenkoerper{data/flow_azn_schwelle}
\bottomrule
\end{tabular}
\end{table}

\begin{table}[htbp]
\centering
\caption{Dieselbe Schwelle unter der zweiten Endwertannahme: die Aktie notiert
am Ende des Horizonts nominal zum heutigen Kurs. Das ist eine Annahme, keine
Prognose. Sie besagt nur, dass es weder Bewertungsgewinn noch
Bewertungsverlust gibt.}
\label{tab:azn-schwelle-variante}
\begin{tabular}{lrrrr}
\toprule
Niveau & je Aktie & 5 Jahre & 10 Jahre & 15 Jahre\\
\midrule
\tabellenkoerper{data/flow_azn_schwellevariante}
\bottomrule
\end{tabular}
\end{table}

\FloatBarrier

\begin{table}[htbp]
\centering
\caption{Umkehrung eins: der Endwert je Aktie, den der heutige Kurs bereits
voraussetzt, in USD und als Vielfaches des Buchwerts}
\label{tab:azn-endwert}
\begin{tabular}{lrrr}
\toprule
Niveau & 5 Jahre & 10 Jahre & 15 Jahre\\
\midrule
\tabellenkoerper{data/flow_azn_endwert}
\bottomrule
\end{tabular}
\end{table}

\begin{table}[htbp]
\centering
\caption{Umkehrung zwei: das j\"ahrliche Wachstum des Mittelzuflusses in
Prozent, das der heutige Kurs voraussetzt}
\label{tab:azn-wachstum}
\begin{tabular}{lrrr}
\toprule
Niveau & 5 Jahre & 10 Jahre & 15 Jahre\\
\midrule
\tabellenkoerper{data/flow_azn_wachstum}
\bottomrule
\end{tabular}
\end{table}

\FloatBarrier

\subsection{Was die f\"unf Tabellen zusammen sagen}

\bk{Die Rangfolge ist \"uber alle Horizonte und alle drei Niveaus stabil: Die
Rendite steigt mit dem Horizont und mit dem Niveau, und sie erreicht die
H\"urde in keiner einzigen Zelle. Das Niveau der Zahlen h\"angt an der
Endwertannahme, die Rangfolge nicht, und das ist der Grund, dem Ergebnis zu
trauen und zugleich seinem Betrag zu misstrauen.}

\bk{Wie stark der Endwert das Ergebnis tr\"agt, zeigt der Vergleich der beiden
Schwellentabellen. Mit dem Buchwert als Endwert liegt die Zehnjahresschwelle
bei \USDje{\AznSchwelleZehn}; mit der Annahme, dass die Aktie in zehn Jahren
noch zum heutigen Kurs notiert, bei \USDje{\AznSchwelleVarZehn}. Beide liegen
unter dem Kurs von \USDje{\AznKurs}, aber die zweite liegt ihm deutlich
n\"aher. Ein Leser, dem die erste Zahl unglaubw\"urdig erscheint, hat einen
Punkt; und die zweite Zahl bleibt trotzdem eine Schwelle unter dem Kurs.}

\bk{Tabelle~\ref{tab:azn-wachstum} ist die Zeile, an der sich das Votum
pr\"ufen l\"asst. Der heutige Kurs setzt \"uber zehn Jahre
\Pct{\AznWachstumZehn} Wachstum des freien Mittelzuflusses pro Jahr voraus.
Das Unternehmen hat zwischen 2019 und 2025 \Pct{\AznCagr} pro Jahr geliefert.
Die geforderte Rate ist also nicht absurd, sondern rund eineinhalbmal so hoch
wie die tats\"achlich erreichte, und die tats\"achlich erreichte stammt aus
einem Zeitraum, der die \"Ubernahme von Alexion und den Aufstieg der Onkologie
enth\"alt und am h\"ochsten Punkt der Reihe endet.}
